\documentclass[10pt,a4paper]{amsart}
\usepackage[margin=19mm]{geometry}
\usepackage{amsmath,amssymb,mathtools}
\usepackage[scr=rsfs,cal=euler]{mathalfa}
\usepackage{xfrac}
\usepackage[unicode,hidelinks,bookmarks=false]{hyperref}
\newcommand{\ie}{i.e.\ }
\newcommand{\cf}{cf.\ }
\newcommand{\resp}{respectively\ }
\newcommand{\Subsection}{\subsection}
\providecommand{\nobreakdash}{\nobreak\hbox{-}}
\setlength{\emergencystretch}{2em}
\multlinegap0pt
\usepackage{bm}
\usepackage{bbm}
\usepackage{dsfont}
\usepackage{xspace}
\usepackage{cancel}
\usepackage{mathtools}
\usepackage{amsthm}
\makeatletter
\@ifundefined{theorem}{\newtheorem{theorem}{Theorem}}{}
\@ifundefined{lemma}{\newtheorem{lemma}[theorem]{Lemma}}{}
\makeatother
\newcommand{\A}{{\mathcal A}}
\newcommand{\Bs}{{\mathscr{B}}}
\newcommand{\GG}{{\mathcal G}}
\newcommand{\Kk}{{\mathbb{K}}}
\title[Brjuno condition through best approximations and the lineari]{Brjuno condition through best approximations and the linearization problem}
\author{Nicolas Chevallier, Jo\~{a}o Lopes Dias, Jos\'e Pedro Gaiv\~ao, Antoine Marnat, Nikolay Moshchevitin}
\date{Source: arXiv:2607.25610v1 (CC BY 4.0)}
\begin{document}
\maketitle

\noindent\emph{Source and licence.} Brjuno condition through best approximations and the linearization problem. Version arXiv:2607.25610v1 (CC BY 4.0), distributed under the Creative Commons Attribution 4.0 International licence, as recorded in the arXiv licence metadata for this exact version and shown on the abstract page https://arxiv.org/abs/2607.25610v1. Full rights notice: licenses/arxiv-2607.25610-CC-BY-4.0.txt.\par\smallskip

\noindent\textit{Editorial scope.} Counted unit: the Lemma whose source label is \texttt{lem:DFu} in the article (source0.tex lines 2548--2554, source label \texttt{lem:DFu}), delivered together with the article's own complete original proof of it (source0.tex lines 2556--2638, one \texttt{proof} environment of 83 lines). No mathematics is written, repaired or re-derived: statement and proof are the article's own text line for line. Required context retained uncounted: lem:DF0 (lines 2464--2468); lem:DFu0 (lines 2501--2506). Every definition, display and label of those lines is retained unchanged. Omitted from this package and not redistributed: 1--2463, 2469--2500, 2507--2547, 2555--2555, 2639--2764. Nothing inside a retained excerpt is omitted or altered: every retained body is delivered exactly as the article's lines are written, so no omission receipt of any kind accompanies this package. Citations: the retained excerpts carry no citation command at all, so no bibliography entry of the article is transcribed and no reference is redistributed. Reference adaptation: none was needed and none was made; every reference inside the delivered unit resolves inside it, so no unresolved reference or citation is produced. Printed slips kept literal: no printed word, symbol, hypothesis or number of the counted statement or of its proof was altered. Layout and class: the article's own class is \texttt{amsart} with packages including amsmath, amssymb, amsthm, amsfonts, graphicx, mathrsfs, color, geometry; the wrapper is the lane's pinned \texttt{amsart} editorial wrapper at 10pt on A4, which restates the theorem-like environments the retained text needs and adds the article's own macro definitions that the retained text needs (\texttt{\textbackslash A}, \texttt{\textbackslash Bs}, \texttt{\textbackslash GG}, \texttt{\textbackslash Kk}), with extra wrapper packages (bm, bbm, dsfont, xspace, cancel, mathtools). The delivered numbering is the unit's own. Assets: no retained span contains a figure environment, a graphics-inclusion command, a PDF-page inclusion, a table or a TikZ picture; no image file is copied into this package, the package declares no asset dependency and no image file is redistributed with it. Licence: version arXiv:2607.25610v1 of this article is distributed under the Creative Commons Attribution 4.0 International licence, as recorded in the arXiv licence metadata for this exact version and shown on the abstract page https://arxiv.org/abs/2607.25610v1. A full rights notice is delivered at licenses/arxiv-2607.25610-CC-BY-4.0.txt. Characters: every retained excerpt is pure ASCII, so the delivered engine, XeLaTeX with its default Latin Modern text font, reports no missing character.
\begin{lemma}\label{lem:DF0} $D\GG(0)^{-1}$ is a bounded linear operator from $\Kk^-_\sigma\A_{\rho}$ to $\Kk^-_\sigma\A'_{\rho}$. Moreover 
$$
|D\GG(0)^{-1}|\leq \frac4\sigma.
$$
\end{lemma}

\begin{lemma}\label{lem:DFu0}
If $u\in\Bs_{\delta}$, then the linear operator $D\GG(u)-D\GG(0)$ mapping $\Kk_\sigma^-\A'_{\rho}$ to $\Kk_\sigma^-\A_{\rho}$ is bounded and
$$
|D\GG(u)-D\GG(0)|\leq \frac{\sigma}{8}.
$$
\end{lemma}

\begin{lemma}\label{lem:DFu}
If $u\in\Bs_{\delta}$, then $D\GG(u)^{-1}$ is a bounded linear operator from $\Kk_\sigma^-\A_{\rho}$ to $\Kk_\sigma^-\A'_{\rho}$ with
\begin{align*}
|D\GG(u)^{-1}|\leq \frac{8}{\sigma}.
\end{align*}
Moreover,  $\Bs_{\delta} \ni u\mapsto D\GG(u)^{-1}$ is locally Lipschitz.
\end{lemma}

\begin{proof}
By Lemmas~\ref{lem:DF0} and \ref{lem:DFu0},
\begin{align*}
| (D\GG(u)-D\GG(0))D\GG(0)^{-1} |
\leq \frac{\sigma}{8}\cdot\frac{4}{\sigma}= \frac12.
\end{align*}

Hence the operator
\[
I+(D\GG(u)-D\GG(0))D\GG(0)^{-1}
\]
is invertible and its inverse is given by the convergent Neumann series
\[
\left[I+(D\GG(u)-D\GG(0))D\GG(0)^{-1}\right]^{-1}
=
\sum_{n\ge 0} \left[-(D\GG(u)-D\GG(0))D\GG(0)^{-1}\right]^n.
\]

Therefore
\begin{align*}
D\GG(u)^{-1}
&= D\GG(0)^{-1}
\left[I+(D\GG(u)-D\GG(0))D\GG(0)^{-1}\right]^{-1}.
\end{align*}

%\medskip

%\noindent
%\textbf{Uniform bound.}
From the Neumann series and the estimate above,
\[
\left|
\left[I+(D\GG(u)-D\GG(0))D\GG(0)^{-1}\right]^{-1}
\right|
\leq \sum_{n\ge 0} (1/2)^n = 2.
\]
Hence
\[
|D\GG(u)^{-1}|
\leq 2|D\GG(0)^{-1}|
\leq \frac{8}{\sigma}.
\]

%\medskip
%
%\noindent
%\textbf{Local Lipschitz dependence.}
Let \(u,v \in \Bs_\delta\). 
%Using the resolvent identity
%\[
%A^{-1}-B^{-1} = A^{-1}(B-A)B^{-1},
%\]
%with \(A=D\GG(u)\), \(B=D\GG(v)\), we obtain
Notice that
\begin{align*}
D\GG(u)^{-1}-D\GG(v)^{-1}
=
D\GG(u)^{-1}\bigl(D\GG(v)-D\GG(u)\bigr)D\GG(v)^{-1}.
\end{align*}

Taking norms and using the uniform bound above,
\[
|D\GG(u)^{-1}-D\GG(v)^{-1}|
\le
\frac{64}{\sigma^2}
|D\GG(u)-D\GG(v)|.
\]

Since \(u \mapsto D\GG(u)\) is continuous and differentiable on \(\Bs_\delta\),
it is locally Lipschitz; hence there exists \(L>0\) such that
\[
|D\GG(u)-D\GG(v)|\le L|u-v|'_\rho.
\]

Combining the estimates yields
\[
|D\GG(u)^{-1}-D\GG(v)^{-1}|
\le \frac{64L}{\sigma^2}|u-v|'_\rho.
\]

Thus \(u \mapsto D\GG(u)^{-1}\) is locally Lipschitz.

\end{proof}

\end{document}
